\scanpage{65}
% archive foliation: 64
\uline{Supposons} que $W^a(W(S)[t])$ existe. C'est alors de façon naturelle un préschéma en anneaux sur $S$, que nous noterons $W$, appelé \uline{préschéma d'anneaux générateur du foncteur} $W$.

[Dans l'exemple II, c'est le schéma d'anneaux $W$ de départ. Dans le cas de l'exemple I, \deleted{c'est le} si \deleted{schéma} on suppose $S'$ \inserted{fini et} loc. libre sur $S$, c'est le \deleted{\unsure{faisc.}} schéma en anneaux $V(\check{\underline A})$, où $\underline A$ est le faisceau d'algèbres sur $S$ associé à $S'$].

Soit alors $\mathfrak X$ ss-schéma fermé de $W(S)[t_1,\ldots,t_r]$ défini par \deleted{\unsure{un}} le faisceau d'idéaux \deleted{\unsure{de}} $\underline O_{W(S)}[t_1,\ldots,t_r]$ engendré par des sections $F_i$ $(i=1,\ldots s)$, donc $\mathfrak X$ est le produit fibré, sur $\mathfrak Y^s$, de $\mathfrak Y^r$ et $\mathfrak Y^0$ [où $\mathfrak Y^{\deleted{\unsure{*}}}=W(S)[t]$].
